% Second pass (Claude, 2026-11-10): first pass checked against the archive images; changes in LM_C3-568-586_changes.md.
\scanpage{568}
% source-position: top, lines 1--2
immédiat sur (77). \deleted{Il reste donc à}\inserted{On a donc}
\deleted{\unsure{d}}

% source-position: upper, display (79)
\[
(79)\qquad
\diagram{\begin{tikzcd}
S_0\Gamma_{\rho,\sigma} \arrow[r,hook] & S\unsure{\mathfrak{G}}_{\rho,\sigma}^{\mathrm{eff}}
\end{tikzcd}}
\]
% source-position: upper, lines 3--5
il reste à voir que cette application est surjective ---
on a fait ce qu'il fallait pour. Notons
d'abord que le composé
\[
\diagram{\begin{tikzcd}
S_0\Gamma_{\rho,\sigma} \arrow[r,hook] & S\unsure{\mathfrak{G}}_{\rho,\sigma}^{\mathrm{eff}} \arrow[r,"\text{épi}"] & \Sigma_{\rho,\sigma}^{\mathrm{eff}}\bigl( \arrow[r,"\sim"] & Z_{\rho,\sigma}\bigr)
\end{tikzcd}}
\]
\deleted{\(\diagram{\begin{tikzcd}(u,a,b)\arrow[r,mapsto]&u\end{tikzcd}}\)}
% source-position: upper-middle, lines 6--7
n'est autre, par (77) et (71) que
\deleted{\(\diagram{\begin{tikzcd}
(\unsure{\alpha},\unsure{\beta},\unsure{\mu})\arrow[r,mapsto]&
(\unsure{[\alpha,\rho]},\unsure{[\beta,\sigma]})
\end{tikzcd}}\)}
\[
(80)\qquad
\diagram{\begin{tikzcd}
x=(u,a,b)\arrow[r,mapsto]& (\xi,\eta,\mu,\varepsilon',\varepsilon)
\end{tikzcd}}
\quad
\left\{\begin{array}{ll}
\xi=[u,\rho]&\text{donc }u(\rho)=\xi\rho\\
\eta=[u,\sigma]&\text{donc }u(\sigma)=\eta\sigma\\
\mu=\chi_{S\Gamma}(x)&\\
\varepsilon'=\varepsilon'(x)&\\
\varepsilon=\varepsilon(x)&
\end{array}\right.
\]
% source-position: middle, lines 8--10
donc son \deleted{image} composé avec
\(\diagram{\begin{tikzcd}\Sigma_{\rho,\sigma}^{\mathrm{eff}}\arrow[r,"\sim"]&Z_{\rho,\sigma}\end{tikzcd}}\)
n'est autre que la projection canonique
\[
\diagram{\begin{tikzcd}
S_0\Gamma_{\rho,\sigma}\arrow[r]&Z_{\rho,\sigma}\\
(u,a,b)\arrow[r,mapsto]&u
\end{tikzcd}}
\]
qui est surjective, de noyau
\(SN_\rho^*\times SN_\sigma^*=SZ_\rho\times SZ_\sigma\).

% source-position: lower-middle, lines 11--12
Considérons d'autre part \deleted{\unsure{l'opération}}
\deleted{Il reste à voir que}
\[
(81)\quad
\left\{\begin{array}{l}
\text{opération de }SN_\rho\times SN_\sigma\\
\text{sur }\deleted{S}\unsure{\mathfrak{G}}_{\rho,\sigma}
\subset\mathcal{G}\times\mathcal{G}\times\widehat{\mathbb Z}^{*}\\
\hfill\Vert\\
\{(\alpha,\beta,\mu)\mid[\alpha,\rho]=[\beta,\sigma]\varepsilon_1^{\mu-1}\}
\end{array}\right.
\quad
\begin{array}{l}
(a_0,b_0)\in S\deleted{\unsure{N}}Z_\rho^{\unsure{*}}\times S\deleted{\unsure{N}}Z_\sigma^{\unsure{*}}\\
\inserted{(=SZ_\rho\times SZ_\sigma)}\\
\text{opère par}\\
\diagram{\begin{tikzcd}(\alpha,\beta,\mu)\arrow[r,mapsto]&(\alpha a_0,\beta b_0,\mu)\end{tikzcd}}
\end{array},
\]
% source-position: lower, lines 13--17
Il est clair qu'il opère \deleted{fidèlement} librement, et
que le quotient de \(\unsure{\mathfrak{G}}_{\rho,\sigma}\) par cette action
s'identifie à \(\Sigma_{\rho,\sigma}\), donc on a le résultat
idem pour \(\unsure{\mathfrak{G}}_{\rho,\sigma}^{\unsure{*}}\) et
\(\Sigma_{\rho,\sigma}^{\unsure{*}}\) \deleted{\illegible{1}}. Or, pour
% source-position: bottom, overwritten display (82), lines 18--20
\deleted{(82) pour \(\varepsilon,\varepsilon'\)
\inserted{\(\varepsilon=(-1)^\nu,\quad\varepsilon'=(-1)^{\nu'}\)} fixés,
l'action (81) de \(SZ'_\rho\times SZ'_\sigma\) correspond
dans \(S_0\Gamma_{\rho,\sigma}\) à la \deleted{translation}\inserted{opération}
\unsure{\((u,a,b)\)}
\(\diagram{\begin{tikzcd}
(u,a,b)\arrow[r,mapsto]&
\bigl(u,\unsure{\varepsilon_\rho^{\tau'}(a_0^{-1})}a,
\unsure{\varepsilon_\sigma^{\tau'}(b_0^{-1})}b\bigr)
\end{tikzcd}}\)}
\mspage{567}

% Printed archive footer outside the manuscript frame (not manuscript):
% 567
% https://grothendieck.umontpellier.fr
% Université de Montpellier, UMR CNRS 5149, Case courrier 051, Place Eugène Bataillon
% 34095 Montpellier cedex 5 - France

\scanpage{569}
% source-position: top edge, marginal calculation
\margin{\deleted{\unsure{i.e.}\ \(\unsure{\varepsilon_\rho^{\tau'}}\,
\unsure{(a)}=\tau'^{\nu'}\)}}
\mspage{588}
% source-position: top line 1, cancelled
\deleted{\unsure{c'est donc} (\illegible{2} \unsure{l'ensemble}}
% source-position: upper display (82)
\[
(82)\qquad \varepsilon=(-1)^\nu,\quad\varepsilon'=(-1)^{\nu'}
\]
fixés, l'opération précédente de \((a_0,b_0)\) correspond,
par l'application (76), à l'opération
\[
(83)\quad
\diagram{\begin{tikzcd}
(u,a,b)\arrow[r]&\bigl(u,\tau^{\nu'}(a_0^{-1})a,\tau^\nu(b_0^{-1})b\bigr)
\end{tikzcd}},
\]
% source-position: upper-middle, lines 5--8
i.e. à la translation à gauche par
\((\tau^{\nu'},\tau^\nu)\underbrace{(a_0^{-1},b_0^{-1})}_{\inserted{(a_0,b_0)^{-1}}}\).
D'où la surjectivité, et l'isomorphisme cherché
\[
(84)\qquad
\diagram{\begin{tikzcd}
S_0\Gamma_{\rho,\sigma}\arrow[r,"\sim"]&S\unsure{\mathfrak{G}}_{\rho,\sigma}^{\mathrm{eff}}
\end{tikzcd}}
\quad\text{(cf. 76)}
\]
% source-position: middle, lines 9--12
La translation à gauche par \(\ell\in L_0\)
\inserted{dans \(S_0\Gamma_{\rho,\sigma}\)} correspond,
par la bijection précédente, à l'\deleted{translation}\inserted{opération}
\[
(85)\qquad
\diagram{\begin{tikzcd}
(\alpha,\beta,\mu)\arrow[r,mapsto]&(\ell\alpha,\ell\beta,\mu)
\end{tikzcd}}
\]
dont il était clair à priori qu'elle stabilise
\(S\unsure{\mathfrak{G}}_{\rho,\sigma}^{\mathrm{eff}}\).
Passant au quotient, on trouve donc
\[
(86)\qquad
\diagram{\begin{tikzcd}
\Gamma_{\rho,\sigma}\arrow[r,"\sim"]&L_0\backslash S\unsure{\mathfrak{G}}_{\rho,\sigma}^{\mathrm{eff}}
\end{tikzcd}}.
\]

% source-position: lower, lines 13--14 and display (87)
Enfin, considérons
\[
(87)\quad
\left\{\begin{array}{l}
\diagram{\begin{tikzcd}[column sep=large]
S_0M_{\rho,\sigma}\arrow[r]&S\unsure{\mathfrak{G}}_{\rho,\sigma}^{\mathrm{eff}}\times\mathcal{G}^+
=\underbrace{\{\overset{\inserted{(\in\mathcal{G}\times\mathcal{G}\times\widehat{\mathbb Z}^{*}\times\mathcal{G}^{+}}}{(\alpha,\beta,\mu,f)}\mid}
_{\unsure{(}[\alpha,\rho]=[\beta,\sigma]\varepsilon_1^{\mu-1}\unsure{)}}\\
x=(u,a,b,U)\arrow[r,mapsto]&
\deleted{((\unsure{ua^{-1}\tau'},\unsure{ub^{-1}\tau'}),\unsure{uU^{-1}})}
\inserted{((\alpha,\beta,\mu),f)}
\end{tikzcd}}
\\
\hfill
\left\{\begin{array}{l}
\alpha=ua^{-1}\tau'^{\deleted{\unsure{\nu'}}\inserted{\unsure{\nu'_\rho(a)}}}\\
\beta=ub^{-1}\tau'^{\unsure{\nu_\sigma(b)}}\\
\mu=\chi_{SM}(x)\quad(=\chi_Z(u)=\chi_\rho(a)=\chi_\sigma(b)=\chi_G(U))\\
f=uU^{-1}
\end{array}\right.
\end{array}\right.
\]
% source-position: bottom line 15
Cette application est bijective (trivial par (84)).
\mspage{568}

% Printed archive footer outside the manuscript frame (not manuscript):
% 568
% https://grothendieck.umontpellier.fr
% Université de Montpellier, UMR CNRS 5149, Case courrier 051, Place Eugène Bataillon
% 34095 Montpellier cedex 5 - France

\scanpage{570}
\mspage{588}
% source-position: top heading
\textcircled{\(7^\circ\)}\quad\uline{Fonctorialités}

Considérons
\[
(88)\qquad
\diagram{\begin{tikzcd}\mathcal{G}\arrow[r,"\varphi"]&\mathcal{G}'\end{tikzcd}}
\qquad\begin{array}{l}\text{hom surjectif de}\\\text{groupes profinis}\end{array}
\]
% source-position: upper lines 3--5
soient \deleted{\(\unsure{\rho,\sigma,\tau'}\)}
\(\rho_{\mathcal{G}'},\sigma_{\mathcal{G}'},\tau'_{\mathcal{G}'},\ldots\)
les images de \(\rho=\rho_{\mathcal{G}},\sigma=\sigma_{\mathcal{G}},
\tau'=\tau'_{\mathcal{G}},\ldots\) dans \(\mathcal{G}'\).

% source-position: upper-middle line 6
\uline{Prop}\quad Conditions équivalentes sur \(\varphi\) :

a) L'application induite par \(\varphi\)
\[
(89)\qquad
\diagram{\begin{tikzcd}
\Sigma_{\rho,\sigma}^{\mathcal{G}}\arrow[r]&\Sigma_{\rho,\sigma}^{\mathcal{G}'}
\end{tikzcd}}
\]
\deleted{\unsure{transforme}}\inserted{applique}
\(\Sigma_{\rho,\sigma}^{\mathcal{G}\,\mathrm{eff}}\)
dans \(\Sigma_{\rho,\sigma}^{\mathcal{G}'\,\mathrm{eff}}\).

% source-position: middle lines 8--11
b) Pour tout \(u\in\operatorname{Aut}\mathcal{G}^+\),
\unsure{qui|que} normalise \(L_0\),
et \unsure{tel|telle} que \(u(\deleted{\rho}L_\rho)\underset{\mathcal{G}^+}{\sim}L_\rho\),
\(u(L_\sigma)\underset{\mathcal{G}^+}{\sim}L_\sigma\),
il existe \(u'\in\operatorname{Aut}(\mathcal{G}'^+)\) qui rend commutatif
\[
(90)\qquad
\diagram{\begin{tikzcd}
\mathcal{G}^+\arrow[r,"\varphi"]\arrow[d,"u"']&\mathcal{G}'^+\arrow[d,"u'","{\unsure{-}}"']\\
\mathcal{G}^+\arrow[r,"\varphi"]&\mathcal{G}'^+
\end{tikzcd}},
\]
en d'autres termes, \(u\) \deleted{\unsure{applique}} \deleted{\unsure{stabilise}}
\inserted{\unsure{respecte}} \(\operatorname{Ker}\varphi\).

% source-position: middle-lower lines 12--15
b') idem, en supposant seulement, à la place que
\(u\) normalise \(L_0\), que
\(u(L_0)\underset{\mathcal{G}^+}{\sim}L_0\)
(plus les autres conditions)

\unsure{c|b}) \(\exists\) hom de groupes
\[
(91)\qquad
\varphi_Z:\quad
\diagram{\begin{tikzcd}
Z_{\rho,\sigma}^{\mathcal{G}}\arrow[r,"{\deleted{\unsure{\varphi_Z}}}"]&Z_{\rho,\sigma}^{\mathcal{G}'}
\end{tikzcd}}
\]
\deleted{\unsure{avec} \(\unsure{u'}\)}
% source-position: lower lines 16--19
compatible avec \(\chi,\varepsilon_Z,\varepsilon'_Z\) et tel que
\(\forall u\in Z_{\rho,\sigma}^{\mathcal{G}}\), si
\(u'\) est son image on ait comm. dans (90).
(\uline{NB} \(\varphi_Z\) \unsure{sera} \deleted{\illegible{1}}
unique \unsure{par|pour} \unsure{\(\zeta\)|?}, et est compatible avec \(i_{L_0,Z}\))
\margin{donc il induit
\(\diagram{\begin{tikzcd}
\gamma_{\rho,\sigma}^{\mathcal{G}}\arrow[r]&\gamma_{\rho,\sigma}^{\mathcal{G}'}
\end{tikzcd}}\)}

% source-position: bottom lines 20--23
c') \(\exists\) hom de groupes
\[
(92)\qquad
\varphi_{\deleted{Z}\inserted{G}}:\quad
\diagram{\begin{tikzcd}
G_{\rho,\sigma}^{\mathcal{G}}\arrow[r]&G_{\rho,\sigma}^{\mathcal{G}'}
\end{tikzcd}}
\]
compatible avec \(\chi_G,\varepsilon_G,\varepsilon'_G\) et satisfaisant à la
condition de compatibilité (90). (NB ce \(\varphi_G\)
\mspage{569}

% Printed archive footer outside the manuscript frame (not manuscript):
% 569
% https://grothendieck.umontpellier.fr
% Université de Montpellier, UMR CNRS 5149, Case courrier 051, Place Eugène Bataillon
% 34095 Montpellier cedex 5 - France

\scanpage{571}
\mspage{589}
% source-position: top lines 1--4
est alors unique, et compatible avec \(i_{\mathcal{G},G}\)
\inserted{\(i_{Z,G}\)}\unsure{;}
\deleted{\unsure{transforme}}\inserted{applique}
\(\deleted{\unsure{N}}\inserted{Z}_\rho^{\mathcal{G}}\) dans
\(\deleted{\unsure{N}}\inserted{Z}_\rho^{\mathcal{G}'}\),
\(\deleted{\unsure{N}}\inserted{Z}_\sigma^{\mathcal{G}}\) dans
\(\deleted{\unsure{N}}\inserted{Z}_\sigma^{\mathcal{G}'}\),
\unsure{qui} se prolonge en
\(\varphi_\rho:\diagram{\begin{tikzcd}N_\rho^{\mathcal{G}}\arrow[r]&N_\rho^{\mathcal{G}'}\end{tikzcd}}\)
induisant
\(\diagram{\begin{tikzcd}N_\rho^{*\mathcal{G}}\arrow[r]&N_\rho^{*\mathcal{G}'}\end{tikzcd}}\),
et se prolonge en
\(\varphi_\sigma:\diagram{\begin{tikzcd}N_\sigma^{\mathcal{G}}\arrow[r]&N_\sigma^{\mathcal{G}'}\end{tikzcd}}\)
induisant
\(\diagram{\begin{tikzcd}N_\sigma^{*\mathcal{G}}\arrow[r]&N_\sigma^{*\mathcal{G}'}\end{tikzcd}}\)\ldots)

% source-position: upper lines 5--7
\deleted{\unsure{il} \illegible{4}}
\deleted{d) \(\exists\) hom de groupes}

On dira alors que \(\varphi\) est ``prolongeable''.

% source-position: upper-middle lines 8--10
\uline{Corollaire}\quad On peut alors prolonger \(\varphi\) en
\[
(93)\qquad\varphi_M:
\diagram{\begin{tikzcd}
S_0G_{\rho,\sigma}^{\mathcal{G}}\arrow[r]&S_0G_{\rho,\sigma}^{\mathcal{G}'}
\end{tikzcd}},
\]
compatible avec \(\Theta_Z,\Theta_\rho,\Theta_\sigma,\Theta_G\)
(compte tenu de \(\varphi_Z,\varphi_\rho,\varphi_\sigma,\varphi_G\)),
et ceci de façon ``compatible avec \unsure{tout}'', notamment
avec les diagrammes (45) (62).

% source-position: middle-lower lines 11--14
\uline{Exemple}\quad Supposons
\(\Sigma_{\rho,\sigma}^{\mathcal{G}'}=\Sigma_{\rho,\sigma}^{\mathcal{G}'\,\mathrm{eff}}\),
alors les conditions de la prop. sont \unsure{automatiquement}
satisfaites (quel que soit l'hom surjectif
\(\diagram{\begin{tikzcd}\mathcal{G}\arrow[r]&\mathcal{G}'\end{tikzcd}}\)).

% source-position: lower lines 15--19
\uline{Rmq}\quad Il se peut qu'on se donne un sous-groupe
\(Z_{\rho,\sigma}^{\unsure{\natural|\sharp}}\subset Z_{\rho,\sigma}^{\mathcal{G}}\),
et qu'on considère \deleted{\unsure{a)}} la condition a) de la prop.
seulement pour des \(u\in Z_{\rho,\sigma}^{\unsure{\natural|\sharp}}\).
Mais à l'aide de \(Z_{\rho,\sigma}^{\unsure{\natural|\sharp}}\), on construit
\margin{Il faut supposer \unsure{en outre} que
\(\tau'_Z\in Z_{\rho,\sigma}^{*}\)}
\[
(94)\quad
G_{\rho,\sigma}^{\unsure{\natural|\sharp}}=
Z_{\rho,\sigma}^{\unsure{\natural|\sharp}}\cdot\mathcal{G}^+,
\quad Z_\rho^{\unsure{\natural|\sharp}}=Z_\rho\cap G_{\rho,\sigma}^{\unsure{\natural|\sharp}},
\quad Z_\sigma^{\unsure{\natural|\sharp}}=Z_\sigma\cap G_{\rho,\sigma}^{\unsure{\natural|\sharp}}
\]
% source-position: bottom lines 20--22
d'où \(N_\rho^{\unsure{\natural|\sharp}}\deleted{_{\unsure{\rho,\sigma}}}\)
et \(N_\rho^{\unsure{\natural|\sharp}*}\deleted{_{\unsure{\rho,\sigma}}}\),
\(N_\sigma^{\unsure{\natural|\sharp}}\) et
\(N_\sigma^{\unsure{\natural|\sharp}*}\),
\(\gamma_{\rho,\sigma}^{\unsure{\natural|\sharp}}\), puis
\(\unsure{S}_0M_{\rho,\sigma}^{\unsure{\natural|\sharp}}\) et
\(M_{\rho,\sigma}^{\unsure{\natural|\sharp}}\).
\deleted{On a une variante de \unsure{la}}
Ceci posé, on a une variante évidente de la
\mspage{570}

% Printed archive footer outside the manuscript frame (not manuscript):
% 570
% https://grothendieck.umontpellier.fr
% Université de Montpellier, UMR CNRS 5149, Case courrier 051, Place Eugène Bataillon
% 34095 Montpellier cedex 5 - France

\scanpage{572}
% source-position: top lines 1--2
proposition, donnant des conditions de \unsure{prolongt|prolongement} de \(\varphi\) en
\(\diagram{\begin{tikzcd}
S_0M_{\rho,\sigma}^{*}\arrow[r]&S_0M_{\rho,\sigma}^{\mathcal{G}'}
\end{tikzcd}}\).

% source-position: upper abandoned block, lines 3--7
\deleted{\uline{Cas universel}
\[
\mathcal{G}=\operatorname{Gl}(2,\mathbb Z)^\wedge,\quad
\rho=\begin{pmatrix}0&1\\-1&1\end{pmatrix},\quad
\sigma=\begin{pmatrix}0&-1\\1&0\end{pmatrix}
\]
\[
\tau'=\begin{pmatrix}0&1\\1&0\end{pmatrix},\quad
\unsure{\{}\ \tau=\begin{pmatrix}-1&0\\0&+1\end{pmatrix};\quad
\varepsilon_0=\begin{pmatrix}1&-1\\0&1\end{pmatrix},
\]
\deleted{On a}
\[
\varepsilon_1=\begin{pmatrix}\unsure{1}&0\\1&1\end{pmatrix}
\deleted{^{\unsure{\varepsilon}}}.
\]
Comme \(\rho\)}
\mspage{571}

% Printed archive footer outside the manuscript frame (not manuscript):
% 571
% https://grothendieck.umontpellier.fr
% Université de Montpellier, UMR CNRS 5149, Case courrier 051, Place Eugène Bataillon
% 34095 Montpellier cedex 5 - France

\scanpage{573}
\mspage{590}
% source: upper portion, heading, lines 1--5, (1)--(2).

$8^\circ)$ \uline{Le cas universel} $\mathcal{G}=\mathrm{Gl}(2,\mathbb{Z})^{\wedge}$.

\unsure{.} Considérons la suite exacte

\diagram{\begin{tikzcd}
\text{(1)}\quad 1 \arrow[r] & \{\pm1\} \arrow[r] & \mathcal{N}_{1,1}^{\sim} \arrow[r] \arrow[d,equal] & \mathcal{M}_{0,3}^{\sim} \arrow[r] & 1 \\
 & & \mathcal{M} & &
\end{tikzcd}}

\deleted{le groupe $\{\pm1\}$}
On a
\[
(2)\qquad \{\pm1\}=\unsure{\mathrm{Centr}|\mathrm{Centre}}(\mathcal{M})=\unsure{\mathrm{Centr}|\mathrm{Centre}}(\widehat{\mathcal{G}}^{+}),
\]

% source: oblique left margin beside (2)--(3), four identifications.
\margin{\unsure{on a pris|on a posé}
\[
\begin{aligned}
\mathcal{G}&=\unsure{\mathcal{T}_{1,1}|\mathcal{G}_{1,1}}\simeq\mathrm{Gl}(2,\mathbb{Z}),\\
\widehat{\mathcal{G}}&=\widehat{\mathcal{T}}_{1,1}\simeq\mathrm{Gl}(2,\mathbb{Z})^{\wedge},\\
\unsure{\widehat{\mathcal{G}}^{+}|\mathcal{G}^{+}}&=\mathcal{T}_{1,1}^{+}=\mathrm{Sl}(2,\mathbb{Z}),\\
\widehat{\mathcal{G}}^{+}&=\mathrm{Sl}(2,\mathbb{Z})^{\wedge}.
\end{aligned}
\]}

% source: middle portion, lines 6--12, (3)--(4).
d'où par (1) une application

\diagram{\begin{tikzcd}
\text{(3)}\quad \mathcal{M}_{0,3}^{\sim} \arrow[r] & \mathrm{Aut}(\mathcal{M}) \arrow[r,"{\inserted{\substack{\text{restriction à}\\ \mathcal{G}^{+\wedge}=\mathrm{Sl}(2,\mathbb{Z})^{\wedge}}}}"] & \mathrm{Aut}(\widehat{\mathcal{G}}^{+\deleted{\unsure{\wedge}}})
\end{tikzcd}}

\deleted{dans le cas}
D'ailleurs, \unsure{tout} \unsure{autom.} de $\widehat{\mathcal{G}}^{+}$ \unsure{conserve}
son centre, donc \unsure{passe} au quotient, d'où

\diagram{\begin{tikzcd}
\text{(4)}\quad \mathrm{Aut}(\widehat{\mathcal{G}}^{+}) \arrow[r] & \mathrm{Aut}(\widehat{\mathcal{G}}^{+}/\!\pm1\simeq\deleted{\unsure{\mathcal{M}}}\mathcal{G}_{0,3}^{+\wedge})
\end{tikzcd}}

% source: lower-middle portion, lines 13--16, triangle (5).
et \unsure{le composé|la composée} de (3) avec (4) est l'hom
canonique:

\diagram{\begin{tikzcd}
\text{(5)}\quad \mathcal{M}_{0,3}^{\sim} \arrow[rr] \arrow[dr] & & \mathrm{Aut}(\widehat{\mathcal{G}}^{+}) \arrow[dl] \\
& \mathrm{Aut}(\widehat{\mathcal{G}}_{0,3}^{+}=\widehat{\mathcal{G}}^{+}/\{\pm1\}) &
\end{tikzcd}}

Cela montre que (3) est injectif, puisque
$\mathcal{M}_{0,3}^{\sim}\longrightarrow\mathrm{Aut}(\widehat{\mathcal{G}}_{0,3}^{+})$
l'est. \deleted{L'application}

% source: bottom portion, lines 17--20, remark continuing on scan574.
\uline{Rmq} L'application
$\mathrm{Aut}(\widehat{\mathcal{G}}^{+})\longrightarrow\mathrm{Aut}(\mathcal{G}_{0,3}^{+\wedge})$
n'est pas injective par contre, \deleted{\unsure{le}} les
éléments du noyau correspondent aux
\unsure{hom}
$\widehat{\mathcal{G}}^{+}\longrightarrow\unsure{\mathrm{Centr}|\mathrm{Centre}}(\widehat{\mathcal{G}}^{+})=\{\pm1\}$,

\mspage{572}

% Printed archive footer outside the manuscript frame (not manuscript):
% 572
% https://grothendieck.umontpellier.fr
% Université de Montpellier, UMR CNRS 5149, Case courrier 051, Place Eugène Bataillon
% 34095 Montpellier cedex 5 - France

\scanpage{574}
% source: upper portion, lines 1--5, (6)--(7), including right marginal map.

ou encore aux hom \unsure{de}

\diagram{\begin{tikzcd}
\text{(6)}\quad \mathcal{G}_{\mathrm{ab}}^{+\wedge} \arrow[r] & \{\pm1\} & \text{ou encore}\quad (\mathcal{G}_{\mathrm{ab}}^{+})_{2} \arrow[r] & \{\pm1\}
\end{tikzcd}}

Or on voit que
\[
\mathcal{G}_{\mathrm{ab}}^{+\wedge}
=\left\{\rho,\sigma\ \middle|\ \begin{gathered}
\rho^{6}=\sigma^{4}=[\rho,\sigma]=1\\
\rho^{3}=\sigma^{2}
\end{gathered}\right\}=\mathbb{Z}/12\mathbb{Z}
\]

\margin{\diagram{\begin{tikzcd}
\rho \arrow[r,mapsto] & 2\bmod12\\
\sigma \arrow[r,mapsto] & 3\bmod12
\end{tikzcd}}}

donc

\diagram{\begin{tikzcd}
& (\mathcal{G}_{\mathrm{ab}}^{+})_{2} \arrow[r,"\sim"] & \mathbb{Z}/2\mathbb{Z}\\
\text{(7)}\quad\left\{\vphantom{\begin{matrix}1\\1\\1\end{matrix}}\right. & \rho \arrow[r,mapsto] & 0\\
& \sigma \arrow[r,mapsto] & 1
\end{tikzcd}}

% source: middle portion, lines 6--9, composite diagram (8).
\deleted{donc à}
\inserted{Cet isomorphisme correspond au caractère}
$\mathcal{G}^{+\wedge}\longrightarrow\{\pm1\}$ qui s'obtient aussi comme
composé

\diagram{\begin{tikzcd}[column sep=large]
\text{(8)}\quad \mathcal{G}^{+\wedge} \arrow[r] \arrow[rr,bend left=27,shorten >=50pt,"{\unsure{\mathrm{sg}|\mathrm{sgn}}}"] &
\underset{\substack{\text{groupe symétrique}\\\text{à trois lettres}}}{\mathcal{G}_{3}}
=\mathcal{G}^{+\wedge}/\underset{\substack{\text{image}\\\text{inverse}\\\text{de }\Pi_{0,3}^{\wedge}\\\text{dans }\mathcal{G}^{+\wedge}}}{\widehat{\Pi}}
\arrow[r,"{\substack{\text{signature}\\\text{\unsure{d'une} perm.}}}"] & \{\pm1\}
\end{tikzcd}}

% source: middle-lower portion, lines 10--13, (9)--(10).
Ainsi le seul élément \inserted{$\Theta$} \deleted{\unsure{différent}} de $1$
du noyau de \deleted{(3) est donné par} \inserted{(4) est donné par}

\diagram{\begin{tikzcd}
\text{(9)}\quad \deleted{\mathrm{Aut}\mathcal{G}}\ \inserted{\Theta:}\ x \arrow[r,mapsto] & x\cdot\mathrm{sg}(x)
\end{tikzcd}}

On notera, \deleted{que \unsure{$\Theta$ et}} en posant
\[
(10)\qquad L_{0}^{\mathcal{G}}=\varepsilon_{0}^{\widehat{\mathbb{Z}}},\qquad
L_{0}^{\Pi}=L_{0}^{\mathcal{G}}\cap\Pi=l_{0}^{\widehat{\mathbb{Z}}},
\]
% source: lower portion, heavily rewritten lines 14--17, (11), ending unfinished.
que le s-groupe $\Theta(L_{0}^{\mathcal{G}})$\inserted{$\subset\mathcal{G}^{+\wedge}$}
n'est \uline{pas} $\mathcal{G}$-conjugué à $L_{0}$,
\deleted{\unsure{ou}} \deleted{et $\Theta(L_{0}^{\Pi})$ $\mathcal{G}$-conjugué}
\inserted{car} \inserted{il est engendré par}
\deleted{\unsure{$\Theta\varepsilon$}}\deleted{et a fortiori pas.} :
\[
(11)\qquad \Theta(\varepsilon_{0})=-\varepsilon_{0}
\qquad\left(\begin{gathered}
\text{car }\varepsilon_{0}=\sigma\rho\\
\mathrm{sg}(\varepsilon_{0})=-1
\end{gathered}\right)
\]
et on vérifie que $\varepsilon_{0}$ n'est pas conjugué à
$(-\varepsilon_{0})^{\mu}=\deleted{\unsure{\varepsilon_{0}}}-\varepsilon_{0}^{\mu}$
pour \unsure{quelque} $\mu\in\widehat{\mathbb{Z}}^{*}$ (comme on voit

\mspage{573}

% Printed archive footer outside the manuscript frame (not manuscript):
% 573
% https://grothendieck.umontpellier.fr
% Université de Montpellier, UMR CNRS 5149, Case courrier 051, Place Eugène Bataillon
% 34095 Montpellier cedex 5 - France

\scanpage{575}
\mspage{591}
% source: upper portion, pale continuation, lines 1--9 and (12).

tout de suite, en \deleted{passant: $\mathcal{G}/\Pi_{\unsure{a}}$}\inserted{passant:}
$\mathcal{G}_{0,3}^{\wedge}/\Pi_{0,3}^{\wedge}{}_{\mathrm{ab}}$
et \unsure{notant} que $l_{0}=\varepsilon_{0}^{2}$ serait
conjugué à: $(-\varepsilon_{0}^{\mu})^{2}=l_{0}^{\mu}$,
\unsure{qu'on doit avoir|qui en doit avoir} $\mu=1$, or
$\varepsilon_{0}$ et $-\varepsilon_{0}$ ne sont pas conjugués
même dans $\mathrm{Sl}(2,\widehat{\mathbb{Z}})$\unsure{'|!} ... ). Or on va
s'intéresser \unsure{précisément} aux autom \unsure{de|$\varphi$}
$\mathcal{G}^{+\wedge}$ qui transforment $L_{0}^{\mathcal{G}}$ en un sous-groupe
conjugué. On \unsure{notera|note} que par contre
\[
(12).\qquad \Theta(l_{0})=l_{0}\qquad\text{donc}\qquad
\Theta(L_{0}^{\Pi})=L_{0}^{\Pi}
\]

% source: middle portion, proposition and conditions a)--b), lines 10--17.
\uline{Proposition} L'hom (3)

\diagram{\begin{tikzcd}
\mathcal{M}_{0,3}^{\sim} \arrow[r] & \mathrm{Aut}(\widehat{\mathcal{G}}^{+})
\end{tikzcd}}

est injectif, et son image est formée des
autom\inserted{$u$} de $\widehat{\mathcal{G}}^{+}$ qui satisfont les deux
conditions suivantes

a) $u$ stabilise $\widehat{\Pi}=\mathcal{G}^{+\wedge}[2]
=\mathrm{Ker}(\mathcal{G}^{+\wedge}\longrightarrow\mathrm{Sl}(2,\mathbb{Z}/2\mathbb{Z}))$

b) $u$ transforme $L_{0}^{\mathcal{G}}$ en un s-groupe de $\mathcal{G}^{+\wedge}$
$\mathcal{G}^{+\wedge}$-conjugué de $L_{0}$; i.e.
$\exists\mu\in\widehat{\mathbb{Z}}^{*}$ et $f_{0}\in\mathcal{G}^{+\wedge}$
tels qu'on ait
\[
(13)\qquad u(\varepsilon_{0})=\mathrm{int}(f_{0}^{-1})(\varepsilon_{0}^{\mu}).
\]

% source: lower portion, proof opening, lines 18--21.
\uline{Dém.} L'injectivité a déjà été vue.
\deleted{\unsure{Prouvons|Rappelons}}\inserted{Montrons}
que les conditions a)b) sont nécessaires pour
que $u$ \deleted{\unsure{appartienne}}\inserted{provienne}
d'un $u_{0}\in\mathcal{M}_{0,3}^{\sim}$,
\deleted{Pour} a) c'est trivial, puisque $u_{0}$ stabilise
$\Pi_{0,3}^{\unsure{\wedge|\sim}}$, donc son

\mspage{574}

% Printed archive footer outside the manuscript frame (not manuscript):
% 574
% https://grothendieck.umontpellier.fr
% Université de Montpellier, UMR CNRS 5149, Case courrier 051, Place Eugène Bataillon
% 34095 Montpellier cedex 5 - France

\scanpage{576}
% source: upper portion, continuation of proof, lines 1--7 and (14).

image inverse dans $\mathcal{G}^{+\wedge}$. Pour la condition
b), on est ramené au cas où $u_{0}\in\mathcal{M}_{0,3}^{\sim}[0]$
(quitte à conjuguer par un élément de
$\mathcal{G}_{0,3}^{+\wedge}\simeq\mathcal{G}^{+\wedge}/\!\pm1$),
dans ce cas on a par construction de $\mathcal{N}_{1,1}^{\sim}$ à partir
de $\mathcal{M}_{0,3}^{\sim}$ (cf. \uline{\(\overline{\mathrm{III}}\)}, \unsure{pp.|p.p} 665...)
\[
(14)\qquad u_{0}(\varepsilon_{0})=\varepsilon_{0}^{\mu}
\]
où $\mu\in\widehat{\mathbb{Z}}^{*}$ est le multiplicateur, OK.

% source: middle portion, converse argument, lines 8--17.
Inversement, soit $u$ un autom de $\mathcal{G}^{+\wedge}$ qui
satisfait a),b), considérons l'autom. $\widetilde{u}$ de
$\mathcal{G}^{+\wedge}/\!\pm1\simeq\mathcal{G}_{0,3}^{+\wedge}$
qu'il induit. On voit donc
que $\widetilde{u}(\Pi_{0,3}^{\wedge})=\Pi_{0,3}^{\wedge}$, et que
$\widetilde{u}(\varepsilon_{0})$ est $\mathcal{G}_{0,3}^{\wedge}$-conjugué
à un $\varepsilon_{0}^{\mu}$ ($\mu\in\unsure{\widehat{\mathbb{Z}}^{*}|\widehat{\mathbb{Z}}}$).
\deleted{Soit} Donc par
définition, on a $\widetilde{u}\in\mathcal{M}_{0,3}^{\sim}$
(cf. \unsure{pp.|p.p} 390 ff).
Soit $u'$ l'autom de $\mathcal{G}^{+\wedge}$ induit par $\widetilde{u}$
(via (3)) -- alors $u'$ et $u$ induisent le $\overline{\mathrm m}$
autom de $\mathcal{G}_{0,3}^{+\wedge}=\mathcal{G}^{+\wedge}/\{\pm1\}$,
donc on aura $u'=u\circ\Theta$, où
$\Theta$\inserted{$=u^{-1}u'$} est un autom qui

% source: lower-middle portion, numbered properties and conclusion, lines 18--21.
$1^\circ)$ induit l'identité sur $\mathcal{G}^{+\wedge}/\{\pm1\}$, et

$2^\circ)$ transforme $L_{0}^{\mathcal{G}}$ en un sous-groupe conjugué

-- on a vu que cela implique $\Theta=1$, cqfd.

% source: bottom portion, abandoned proposition and exact sequence (15).
\deleted{\uline{Proposition} Soit $\widehat{\Pi}_{0}$ le sous-groupe}
Considérons maintenant la \deleted{\unsure{lemme}}\inserted{suite}
exacte

\diagram{\begin{tikzcd}
\text{(15)}\quad 1 \arrow[r] & \pm1 \arrow[r] & \widehat{\Pi} \arrow[r] \arrow[d,equal] & \Pi_{0,3}^{\wedge} \arrow[r] & 0\\
& & \mathcal{G}^{+\wedge}[2] & &
\end{tikzcd}}

\mspage{575}

% Printed archive footer outside the manuscript frame (not manuscript):
% 575
% https://grothendieck.umontpellier.fr
% Université de Montpellier, UMR CNRS 5149, Case courrier 051, Place Eugène Bataillon
% 34095 Montpellier cedex 5 - France

\scanpage{577}
\mspage{592}
% source: upper portion, lines 1--6, including (16)--(17).
et considérons le scindage donné par

\diagram{\begin{tikzcd}[column sep=large,row sep=large]
(16)\quad\left\{\vphantom{\begin{matrix}\Pi_{0,3}\\l_i\end{matrix}}\right.
& \Pi_{0,3} \arrow[r]
& \Pi=\mathcal{G}^{+}[2]=\operatorname{Ker}\bigl(\mathrm{Sl}(2,\mathbb Z)\longrightarrow\mathrm{Sl}(2,\mathbb Z/2\mathbb Z)\bigr) \\
& l_i \arrow[r,mapsto]
& \deleted{\unsure{\lambda_i|l_i}\overset{\unsure{\mathrm{dfn}}}{=}}\;\lambda_i\overset{\mathrm{dfn}}{=}-l_i
\end{tikzcd}}

Il \deleted{\unsure{est}} y a un tel hom., en vertu de
\[
(17)\qquad
\left\{\begin{aligned}
(-l_\infty)(-l_1)(-l_0)&=1\quad\text{dans }\mathrm{Sl}(2,\mathbb Z)\quad\text{i.e.}\\
l_\infty l_1l_0&=-1
\end{aligned}\right.
\]
\margin{$)$}

et il est évident que c'est un scindage. On a encore de $\overline{\mathrm m}$ un scindage pour $\widehat\Pi$ sur $\widehat\Pi_{0,3}$.

% source: middle portion, lines 7--12, definition (18).
On désigne par $\Pi_0$ l'image de $16)$, i.e.
\[
(18)\qquad
\left\{\begin{aligned}
\Pi_0&=\text{s-groupe de }\Pi\inserted{=\mathcal{G}^{+}[2]}\text{ engendré par}\\
&\lambda_0,\lambda_1,\lambda_\infty\qquad(\text{où }\lambda_i=-l_i)
\end{aligned}\right.
\]
et on a donc un sous-groupe $\widehat\Pi_0$ de $\widehat\Pi$.

% source: middle-lower portion, proposition a)--b), lines 13--18.
\uline{Proposition}\quad a) Le sous-groupe $\Pi_0$ de $\Pi$ est \deleted{distingué}\inserted{d'indice} $2$, et il est distingué dans $\mathcal{G}^{\unsure{\wedge|*}}$.

b) Le sous-groupe $\widehat\Pi_0$ de $\widehat\Pi$ est \deleted{distingué, \unsure{et}}\inserted{d'indice $2$}, et il est distingué dans $\mathcal M$ -- ou, ce qui revient au $\overline{\mathrm m}$, stable sous l'action de $\mathcal M/\pm1\simeq\mathcal M_{0,3}^{\sim}$.

% source: lower portion, proof and (19); right-hand additions at the same height.
\uline{Dém.}\quad Comme
\[
(19)\qquad
\left\{\begin{aligned}
\rho(\lambda_i)&=\lambda_{i+1}\\
\sigma(\lambda_0)&=\lambda_{1}^{\deleted{\unsure{-1}}},\quad
\sigma(\lambda_1)=\lambda_0^{\deleted{\unsure{-1}}}\\
\tau(\lambda_0)&=\lambda_0^{-1},\quad\tau(\lambda_1)=\lambda_1^{-1}
\end{aligned}\right.
\]
\margin{$\begin{aligned}
\sigma(\lambda_\infty)&=\sigma(\lambda_1\lambda_0)^{-1}=\lambda_1^{-1}\lambda_0^{-1}\\
\deleted{\unsure{\tau'}(\lambda_\infty)}&\deleted{=\lambda_\infty^{-1}}\\
\tau(\lambda_\infty)&=\tau(\lambda_1\lambda_0)^{-1}=\lambda_0\lambda_1
\end{aligned}$}

on voit que $\Pi_0$ est \deleted{\unsure{stable}} normalisé par $\rho,\sigma,\tau$, donc par $\mathcal G$. Qu'il soit d'indice $2$ est évident. Donc $\widehat\Pi_0$ est d'indice $2$ dans $\widehat\Pi$ et normalisé par $\widehat{\mathcal G}$. Pour voir qu'il est
\mspage{576}

% Printed archive footer outside the manuscript frame (not manuscript):
% 576
% https://grothendieck.umontpellier.fr
% Université de Montpellier, UMR CNRS 5149, Case courrier 051, Place Eugène Bataillon
% 34095 Montpellier cedex 5 - France

\scanpage{578}

% source-position: upper portion, continuation and display (20).
normalisé par $\mathcal M$ \unsure{tout} entier, comme
\[
\mathcal M=\mathcal M(0)\mathbin{.}\mathcal G^{+\wedge},
\]
on est ramené à vérifier qu'il est normalisé par
$\mathcal M(0)\simeq\mathcal M_{0,3}^{\sim}(0)$.
Or on \unsure{a}, pour $u\in\mathcal M_{0,3}^{\sim}(0)$
\[
(20)\qquad
\left\{
\begin{aligned}
u(\rho)&=\operatorname{int}(\alpha)\rho\\
u(\sigma)&=\varepsilon_2\operatorname{int}(\beta)\sigma\\
u(\varepsilon_0)&=\varepsilon_0^\mu
\end{aligned}
\right.
\qquad
\left\{
\begin{aligned}
\varepsilon_2&=\varepsilon_2(\mu)\\
\alpha,\deleted{\unsure{\beta}}&\in\widehat{\Pi}\mathbin{.}\{1,\tau\}\\
\beta&\in\widehat{\Pi}
\end{aligned}
\right.
\]
% source-position: upper-middle portion, rewritten display (21).
\uline{d'où}
\[
\deleted{u(\ell_0)=}\quad
\deleted{\unsure{(20)}}\quad
(21)\qquad
\left\{
\begin{aligned}
u(\ell_0)&=\ell_0^\mu\\
u(-\ell_0)&=(-\ell_0)^\mu
\quad\text{i.e.}\quad
u(\lambda_0)=\lambda_0^\mu
\quad(\text{car }(-1)^\mu=-1)
\end{aligned}
\right.
\]
\margin{$\unsure{)}$}

% source-position: middle portion, cancelled proof openings and lambda_1 calculation.
\deleted{Il \unsure{reste} à \unsure{voir}}
\deleted{\unsure{En effet|Enfin}}
donc $u(\lambda_0)\in\widehat{\Pi}_0$, il reste à prouver que
$u(\lambda_1)\in\widehat{\Pi}_0$, \unsure{or}
\[
u(\lambda_1)
=u(\rho(\lambda_0))
=\underbrace{(u\rho)}_{\inserted{\deleted{\unsure{u}}(\operatorname{int}(\alpha)\rho)u}}(\lambda_0)
=\operatorname{int}(\alpha)(\rho)
\left(\underbrace{u(\lambda_0)}_{\inserted{\lambda_0^\mu}}\right)
=\deleted{\unsure{\operatorname{int}(\alpha)(\lambda_1^\mu)}}.
\]
\deleted{et on a \unsure{vu} que}
et comme $\widehat{\Pi}_0$ est normalisé par $\mathcal G^\wedge$
et en particulier par $\deleted{\unsure{\rho}}\alpha\rho\alpha^{-1}$,
on trouve bien $u(\lambda_1)\in\Pi_0$.

% source-position: lower portion, corollary (22).
\uline{Corollaire} On a
\[
(22)\qquad
\left\{
\begin{aligned}
\Pi&\simeq\Pi_0\times\mu\\
\widehat{\Pi}&\simeq\widehat{\Pi}_0\times\mu
\end{aligned}
\right.
\]
où $\mu=\{\pm1\}=\operatorname{Centr}(\mathcal G^{+\wedge})$.

% source-position: bottom portion, cancelled and inserted proof of corollary.
\deleted{Il suffit de \unsure{remarquer} que l'élément $\unsure{-1}$
\unsure{de} $\mu$ \unsure{n'est pas} dans $\Pi_0$.}
\inserted{C'est trivial, puisqu'on a un scindage de l'extension centrale (15).}

\mspage{577}

% Printed archive footer outside the manuscript frame (not manuscript):
% 577
% https://grothendieck.umontpellier.fr
% Université de Montpellier, UMR CNRS 5149, Case courrier 051, Place Eugène Bataillon
% 34095 Montpellier cedex 5 - France

\scanpage{579}
\mspage{593}

% source-position: upper portion, introductory word and definitions (23)--(24).
\unsure{Soient|Soit}
\[
(23)\qquad
\left\{
\begin{aligned}
\mu_\tau&=\{1,\tau\}\subset\mathcal G
&D_\tau&=\mu_\tau\times\mu\subset\mathcal G\\
\mu_\sigma&=\{1,\unsure{\sigma|\tau'}\}\subset\mathcal G
&D_\sigma&=\mu_\sigma\times\mu\subset\mathcal G
\end{aligned}
\right.
\]
de sorte que l'on a
\[
(24)\qquad
\begin{aligned}
D_\tau\cap\unsure{\mathcal G^+}
&=D_\tau\cap\Pi=\mu
&D_\tau\cap\Pi_0&=\{1\}\\
D_\sigma\cap\unsure{\mathcal G^+}
&=D_\sigma\cap\Pi=\mu
&D_\sigma\cap\Pi_0&=\{1\},
\end{aligned}
\]
\margin{\unsure{idem} en remplaçant $\mathcal G,\Pi$ par
$\mathcal G^\wedge,\widehat{\Pi}$}
% source-position: upper-middle portion, definitions (25)--(26).
On pose
\[
(25)\qquad
\left\{
\begin{aligned}
\unsure{\Pi_\tau}&=\mu_\tau\mathbin{.}\Pi=D_\tau\mathbin{.}\Pi_0,
&\Pi_{0\tau}&=\mu_\tau\mathbin{.}\Pi_0\\
\Pi_\sigma&=\mu_\sigma\mathbin{.}\Pi=D_\sigma\mathbin{.}\Pi_0,
&\Pi_{0,\sigma}&=\mu_\sigma\mathbin{.}\Pi_0
\end{aligned}
\right.
\]
de sorte que l'on a
\[
(26)\qquad
\left\{
\begin{aligned}
\Pi_\tau&=\Pi_{0\tau}\times\mu\\
\Pi_\sigma&=\Pi_{0,\sigma}\times\mu
\end{aligned}
\right.,
\]
et idem en mettant des $\wedge$.

% source-position: middle-lower portion, normalization and (27).
\uline{Normalisation de $\alpha$}, $\beta$\quad
Les relations $20$, ne déterminent $\alpha,\beta$ que "modulo signe".
On lève cette ambiguïté en prenant
\[
(27)\qquad
\alpha\in\Pi_{0,\tau},\qquad\beta\in\Pi_0.
\]

% source-position: lower portion, quotient presentation (28).
\uline{Le groupe} $\mathcal G/\Pi_0$.
Il est décrit par gén. $\rho,\sigma,\tau^{\unsure{\prime}}$ et relations \unsure{par}
\[
(28)\qquad
\left\{
\begin{aligned}
&\rho^6=\sigma^4=\tau^2=1,\quad \rho^3=\sigma^2\\
&\tau(\rho)=\sigma(\rho^{-1}),\quad \tau(\sigma)=\sigma^{-1}\\
&\tau(\rho)=\rho
\end{aligned}
\right.
\]
% source-position: bottom portion, calculation (29) and revised concluding relations.
Car on a
\[
(29)\qquad
\tau(\rho)=(-\ell_1^{-1})\rho
\]
donc la dernière relation $28$ s'écrit
$\deleted{\tau(\rho)=\rho}\inserted{-\ell_1^{-1}=1}$
i.e. $\lambda_1=1$
(et implique
$\deleted{\unsure{-\ell_\infty=-\ell_0}}\inserted{\lambda_\infty=\lambda_0}=1$
en appliquant $\rho$).

\mspage{578}

% Printed archive footer outside the manuscript frame (not manuscript):
% 578
% https://grothendieck.umontpellier.fr
% Université de Montpellier, UMR CNRS 5149, Case courrier 051, Place Eugène Bataillon
% 34095 Montpellier cedex 5 - France

\scanpage{580}

% source-position: upper portion, presentations (30)--(31).
En termes de $\rho,\sigma,\tau'$ on trouve de $\overline{\mathrm m}$
\[
(30)\qquad
\left\{
\begin{aligned}
&\rho^6=\sigma^4=\tau'^2=1,\quad \rho^3=\sigma^2\\
&\tau'(\rho)=\rho^{-1},\quad\tau'(\sigma)=\sigma^{-1}\\
&\sigma(\rho)=\rho^{-1}
\end{aligned}
\right.
\]
Le sous-groupe $\mathcal G^+/\Pi_0$ a la présentation en $\rho,\sigma$
\[
(31)\qquad
\left\{
\begin{aligned}
&\rho^6=\sigma^4=1,\quad \rho^3=\sigma^2\\
&\sigma(\rho)=\rho^{-1}
\end{aligned}
\right\}
\]

% source-position: upper-middle portion, action question and extension diagram (32).
Quelle est l'opération de $\mathcal M_{0,3}^{\sim}\simeq\mathcal M/\pm1$
sur $\mathcal G^+/\Pi_0$ ?
\deleted{\unsure{c'est une} opération}
\inserted{On a une structure} d'extension
\[
(32)\qquad
\diagram{
\begin{tikzcd}
1 \arrow[r] & \mu \arrow[r] \arrow[d,equal]
& \mathcal G^+/\Pi_0 \arrow[r]
& \mathcal G^+/\Pi \arrow[r] \arrow[d,equal]
& 1\\
& \pm1 && \mathcal G_3 &
\end{tikzcd}
}
\]
% source-position: middle portion, prose following (32).
et les \unsure{autom.} en question conservent cette structure d'extension.
Les \unsure{élts|élmts} de
$\mathcal M_{0,3}^{!\sim}\simeq\mathcal M^!/\{\pm1\}$
opèrent en induisant l'identité dans le quotient $\mathcal G_3$,
et bien \unsure{sûr} aussi dans $\mu$ --- or un tel autom. de
l'extension correspond à un hom
% source-position: lower-middle portion, character maps (33).
\[
(33)\qquad
\diagram{
\begin{tikzcd}
\mathcal G_3 \arrow[r] & \mu \arrow[d,equal]\\
& \pm1
\end{tikzcd}
}
\quad\text{i.e.}\quad
\diagram{
\begin{tikzcd}
\mathcal G_{3\,\mathrm{ab}} \arrow[r] & \mu
\end{tikzcd}
}
\]
% source-position: lower portion, sign character and conclusion (34).
or
\[
\diagram{
\begin{tikzcd}
\mathcal G_{3\,\mathrm{ab}}\arrow[r,"\uline{\mathrm{sg}}"]
&\{\pm1\}=\mu
\end{tikzcd}
}
\]
donc le groupe des autom. de $\mathcal G^+/\Pi_0$ induisant
l'identité dans $\mathcal G_3$ est $\simeq\mathbb Z/2\mathbb Z$,
un générateur \inserted{$\Theta$} est donné par
\[
(34)\qquad
\Theta(x)=\operatorname{sg}(x)\mathbin{.}x
\quad\text{donc}\quad
\left\{
\begin{aligned}
\Theta(\rho)&=\rho\\
\Theta(\sigma)&=-\sigma
\end{aligned}
\right.
\]

\mspage{579}

% Printed archive footer outside the manuscript frame (not manuscript):
% 579
% https://grothendieck.umontpellier.fr
% Université de Montpellier, UMR CNRS 5149, Case courrier 051, Place Eugène Bataillon
% 34095 Montpellier cedex 5 - France

\scanpage{581}
\mspage{594}

% source-position: upper portion, displays (35) and (35 bis).
Comparant avec les formules $(20)$, on trouve
\[
(35)\qquad
\deleted{\unsure{u}}
\left.u\right|_{\mathcal G^{+\wedge}/\widehat{\Pi}_0}
=\Theta^{\varepsilon_2(u)}:
\quad
\diagram{
\begin{tikzcd}
\dot\rho\arrow[r,mapsto]&\dot\rho\\
\dot\sigma\arrow[r,mapsto]&\varepsilon_2(u)\dot\sigma
\end{tikzcd}
}
\]
i.e.
\[
(35\text{ bis})\qquad
\begin{gathered}
[\dot u,\dot\sigma]\\[-0.5ex]
\unsure{\Vert}\\[-0.5ex]
\inserted{\dot u(\dot\sigma)\dot\sigma^{-1}}
\end{gathered}
=\varepsilon_2(u)^{\cdot}
\quad\text{dans }\mathcal G/\Pi_0.
\]
% source-position: upper left slanted margin, commutator formula and warning.
\margin{Cela s'écrit aussi
$[\dot u,\dot x]=\operatorname{sg}(x)^{\unsure{\nu_2(u)}}$
dans $\mathcal G^{+\wedge}/\Pi_0$
$(u\in\mathcal M^!,x\in\mathcal G^{+\wedge})$
(attention, ça \unsure{foire|fait} \unsure{si on y suppose}
\unsure{pour} $u\in\mathcal M^!$ seulement $u\in\mathcal M\ldots$)}

% source-position: upper-middle portion, revised heading and definition (36).
\deleted{Je voudrais}
Reprenons la définition
\[
(36)\qquad
Z=Z_{\rho,\sigma}\stackrel{\mathrm{dfn}}=
\left\{
u\in\operatorname{Aut}(\mathcal G^{+\wedge})\ \left|\
\begin{aligned}
&u(\unsure{L_\rho})\
\inserted{\mathcal G^{+\wedge}\text{-}}\text{conjugué à }\
\deleted{\unsure{\ell_\rho}}L_\rho\\
&u(\unsure{L_\sigma})\
\inserted{\mathcal G^{+\wedge}}\text{-conj. : }L_\sigma\\
&u(\unsure{L_0})=L_0
\end{aligned}
\right.
\right\}
\]
de sorte qu'on aura, pour $u\in Z$
% source-position: middle portion, display (37) and parameter conditions.
\[
(37)\qquad
\left\{
\begin{aligned}
&u(\rho)\quad\mathcal G\text{-conjugué à }\rho^{\varepsilon'}\\
&u(\sigma)\quad\mathcal G\text{-conjugué à }\sigma^{\unsure{\varepsilon}}
\quad(=\varepsilon\sigma)\\
&u(\varepsilon_0)=\varepsilon_0^\mu
\end{aligned}
\right.
\]
avec $\mu\in\widehat{\mathbb Z}^{*}$ et
$\varepsilon,\varepsilon'\in\{\pm1\}$ bien déterminés,
(\unsure{car}
\diagram{
\begin{tikzcd}
\widehat{\mathbb Z}\arrow[r]&\mathcal G\\
n\arrow[r,mapsto]&\varepsilon_0^n
\end{tikzcd}
}
est injectif, $\rho$ et $\rho^{-1}$ ne sont pas
$\mathcal G^{+\wedge}$-conjugués, ni $\sigma$ et $\sigma^{-1}$).
% source-position: lower-middle portion, proof of stabilization and l_0, l_1 calculations.
Je dis qu'un tel $u$ provient de $\mathcal M_{0,3}^{\sim}$ ---
il reste à vérifier, en vertu de la prop..., que $u$ stabilise
$\widehat{\Pi}$, qui est engendré par $\ell_0,\ell_1,\ell_\infty$.
Or on aura
\[
u(\ell_0)=\ell_0^\mu\in\widehat{\Pi}
\]
et
\[
u(\ell_1)=
\underbrace{u\rho}_{\inserted{
\begin{gathered}
\deleted{\unsure{\operatorname{int}(f)\rho u}}\\
\deleted{\unsure{\operatorname{int}(f)\rho^{\varepsilon'}u}}\\
\operatorname{int}(f)(\rho^{\varepsilon'})u\\
f\in\mathcal G^{+\wedge}
\end{gathered}
}}(\ell_0)
=\operatorname{int}(f)(\rho^{\varepsilon'})(u(\ell_0))
=\operatorname{int}(f)(\rho)^{\varepsilon'}(\ell_0^\mu)
\]
or $\mathcal G^\wedge$ et a fortiori
$\operatorname{int}(f)(\rho)^{\varepsilon'}$ normalise
$\deleted{\unsure{\mathcal G}}\widehat{\Pi}$,
donc \deleted{et} $u(\ell_1)\in\widehat{\Pi}$
% source-position: lower portion, l_infinity calculation and final cancelled words.
et comme $u(-1)=-1$,
\inserted{d'où:}
\[
u(\ell_\infty)=-u(\ell_1\ell_0)^{-1}
=-u(\ell_0^{-1})u(\ell_1)^{-1}
\]
on a aussi $u(\ell_\infty)\in\Pi$, OK.

Si on identifie $u$ à un élément de $\mathcal M_{0,3}^{\sim}$,
on voit que
\deleted{\unsure{c'est}} en fait
\deleted{un élément de $\unsure{Z|\mathcal M}$.}

\mspage{580}

% Printed archive footer outside the manuscript frame (not manuscript):
% 580
% https://grothendieck.umontpellier.fr
% Université de Montpellier, UMR CNRS 5149, Case courrier 051, Place Eugène Bataillon
% 34095 Montpellier cedex 5 - France

\scanpage{582}

% source-position: upper portion, continuation from preceding page.
\[
u\in\mathcal M_{0,3}^{\sim}[0]\simeq\mathcal M^![0]\subset\mathcal M
\]
% source-position: large slanted upper-left margin, correction and later-page reference.
\margin{Attention
on n'avait pas introduit précédemment
$\unsure{\mathcal M_{0,3}^{\sim}|\mathcal M_{0,3}}[0]$, mais seulement
$\mathcal M_{0,3}^{!\sim}[0]\subset\mathcal M_{0,3}^{!\sim}$.
Or il \unsure{n'est pas}
vrai que l'on aura ici $u\in\mathcal M_{0,3}^{!\sim}$,
p. ex. \unsure{on peut avoir} $u=\varepsilon_0$.
L'image de $u$ dans $\mathcal G_3$ est dans
\deleted{\unsure{\operatorname{sg}}}
$\{1,\dot\sigma_0\}$, donc en corrigeant
\inserted{par $\unsure{\varepsilon_0}$},
on aura un élément de $\mathcal M_{0,3}^{!\sim}[0]$.
\deleted{On doit}
Cf rectification plus détaillée à p. $609$ ff.
$Z_{\rho,\sigma}^!=\mathcal M^![0]=Z_{\rho,\sigma}\cap\mathcal M^!$.}

% source-position: upper portion, display (38) and its cancelled side conditions.
\deleted{\unsure{on a}} donc $(20)$, de façon plus précise
\[
(38)\qquad
\left\{
\begin{aligned}
&u(\rho)=\operatorname{int}(\alpha)\rho\\
&u(\sigma)=\varepsilon_2\operatorname{int}(\beta)\sigma\\
&u(\varepsilon_0)=\varepsilon_0^\mu
\end{aligned}
\right.
\qquad
\begin{aligned}
&\unsure{\text{avec}}\ \deleted{\unsure{\alpha}}
\quad\alpha\in\widehat{\Pi}_{0\tau},\ \beta\in\widehat{\Pi}_0\\
&\deleted{\unsure{\varepsilon_2}}\quad
\left\{\begin{aligned}
&\mu\in\widehat{\mathbb Z}^{*}\\
&\deleted{\varepsilon_2=\varepsilon_2(\mu),\
\varepsilon_\tau(\alpha)=\varepsilon_3(\mu)}
\end{aligned}\right.
\end{aligned}
\]
% source-position: upper-middle portion, uniqueness paragraph after (38).
\unsure{à} leur tour \unsure{impliquent|impliquant} $(37)$. En fait
les formules \inserted{$(38)$} pour $u(\rho),u(\sigma),u(\varepsilon_0)$
déterminent de façon unique
$\mu\in\widehat{\mathbb Z}^{*}$, $\varepsilon_2\in\{\pm1\}$,
$\alpha\in\widehat{\Pi}_{0,\tau}$, $\beta\in\widehat{\Pi}_0$,
et on aura nécessairement alors
% source-position: middle portion, display (39), exact sequence (40).
\[
(39)\qquad
\varepsilon_2=\varepsilon_2(\mu),\qquad
\varepsilon_\tau(\alpha)=\varepsilon_3(\mu)
\]
où $\varepsilon_\tau$ est l'hom qui intervient dans la suite
\inserted{exacte}
\[
(40)\qquad
\deleted{\unsure{\varepsilon_0}}
\diagram{
\begin{tikzcd}
1\arrow[r]&\widehat{\Pi}_0\arrow[r]&
\widehat{\Pi}_{0,\tau}\arrow[r,"\varepsilon_\tau"]&
\{\pm1\}\arrow[r]&1
\end{tikzcd}
}.
\]

% source-position: lower-middle portion, proposition and component identifications.
Ainsi on a prouvé

\uline{Proposition}\quad
Dans le cas universel $(\mathcal G=\mathcal G^{+\wedge})$
on a
\[
\diagram{
\begin{tikzcd}
Z_{\rho,\sigma}^!&
\mathcal M^![0]\stackrel{\mathrm{dfn}}=
\mathcal M_{0,3}^{!\sim}[0]\arrow[l,"\sim"']
\end{tikzcd}
},
\]
les composants $u,\mu,\varepsilon',\varepsilon$ d'un élément de
$Z_{\rho,\sigma}^!$ correspondent resp.
(en termes de $u_0\in\mathcal M[0]$):
$u_0$ identifié à un élt de
$\operatorname{Aut}(\mathcal G^{+\wedge})$,
$\mu=\chi(u_0)$ (le multiplicateur de $u_0$),
$\varepsilon=\varepsilon_3(\mu)$,
$\varepsilon'=\varepsilon_2(\mu)$.

% source-position: lower portion, corollary and chain of quotients.
\uline{Corollaire}\quad On a
\[
G_{\rho,\sigma}\simeq\mathcal M=\mathcal N_{1,1}^{\sim},
\]
\[
\gamma_{\rho,\sigma}\simeq
\mathcal M/\mathcal G^{+\wedge}
\simeq\mathcal N_{1,1}^{\sim}/\mathcal G^{+\wedge}
\simeq\mathcal M_{0,3}^{\sim}/\mathcal G_{0,3}^{+\wedge}
\simeq\Gamma_{\mathbb Q}^{\sim}.
\]

\mspage{581}

% Printed archive footer outside the manuscript frame (not manuscript):
% 581
% https://grothendieck.umontpellier.fr
% Université de Montpellier, UMR CNRS 5149, Case courrier 051, Place Eugène Bataillon
% 34095 Montpellier cedex 5 - France

\scanpage{583}
\mspage{595}
% source-position: upper heading, Remarque and (41)-(42).

\uline{Remarque} Il est tentant d'écrire les \inserted{deux premières}
relations (38) sous la forme
\[
(41)\qquad
\left\{\begin{array}{ll}
(\alpha^{-1}u)(\rho)=\rho
  & \text{i.e.}\quad \alpha^{-1}u\in Z_\rho^{\unsure{\sharp}}\\
(\beta^{-1}u)(\sigma)=\varepsilon_2\sigma
  & \unsure{\text{d'où}|\text{i.e.}}\quad \beta^{-1}u\in N'_\sigma\\
\deleted{\unsure{u(\varepsilon_0)}} &
\end{array}\right.
\]
et d'introduire alors
\[
(42)\qquad
\left\{\begin{array}{ll}
a=\alpha^{-1}u\in Z_\rho^{\unsure{\sharp}},
  & \text{donc }u=\alpha a,\quad \alpha=ua^{-1}\\
b=\beta^{-1}u\in N'_\sigma,
  & \text{donc }u=\beta b,\quad \beta=ub^{-1},
\end{array}\right.
\]
% source-position: middle paragraph and quotient diagram preceding (43).
et les développements du \S\,48 (notamment
\uline{\(\overline{\mathrm{XI}}\)}) peuvent faire s'attendre que les
quantités \(a,b\), pour \(u\in\mathcal{M}[0]\) \unsure{variable},
dépendent \uline{multiplicativement} de \(u\). Nous allons voir cependant
qu'il n'en est rien. Soit en effet
\diagram{\begin{tikzcd}[column sep=large]
\delta:\ \deleted{\unsure{\mathcal{M}}}\ G_{\rho,\sigma}
  \arrow[r,"\sim"] & \mathcal{M}
  \arrow[r] & \gamma_{\rho,\sigma}\simeq\Gamma_{\mathbb Q}
\end{tikzcd}}
l'hom canonique, de sorte que l'on aura
\[
(43)\qquad
\delta(u)=\delta(b)=\tau_\gamma^{\unsure{\nu_3}}\delta(a),
\]
en posant
\[
\tau_\gamma=\delta(\tau),\qquad
\varepsilon_3=(-1)^{\nu_3}\quad
(\nu_3\in\mathbb Z/2\mathbb Z)
\]
% source-position: lower half, (43 bis), (43 ter), and (44).
Si \(u'\) correspond à \(b',a',\) \inserted{\(\nu'_3\)}, on aura de
\(\overline{\mathrm m}\)
\[
(43\,\mathrm{bis})\qquad
\delta(u')=\delta(b')=\tau_\gamma^{\nu'_3}\delta(a')
\]
et \unsure{si} \(u''=u'u\) correspond à \(a'',b'',\nu''_3\),
\[
(43\,\mathrm{ter})\qquad
\delta(u'u)=\delta(b'')=\tau_\gamma^{\nu''_3}\delta(a''),
\]
En fait, on aura
\[
\begin{array}{rl}
(44)\qquad & b''=b'b,\quad \nu''_2=\nu'_2+\nu_2\quad
 \unsure{\text{et.}|\text{et}}\quad \nu''_3=\nu'_3+\nu_3\\
& (\text{i.e. }\varepsilon''_2=\varepsilon'_2\varepsilon_2
\quad\text{et}\quad\varepsilon''_3=\varepsilon'_3\varepsilon_3)
\end{array}
\]
\mspage{582}

% Printed archive footer outside the manuscript frame (not manuscript):
% 582
% https://grothendieck.umontpellier.fr
% Université de Montpellier, UMR CNRS 5149, Case courrier 051, Place Eugène Bataillon
% 34095 Montpellier cedex 5 - France

\scanpage{584}
% source-position: upper calculations leading to (45).

On a vu
\[
\begin{array}{l}
\deleted{\unsure{\varepsilon_2|\text{et}}}u(\sigma)
 =\varepsilon_2\operatorname{int}(\beta)\sigma\\
u'(\sigma)=\varepsilon'_2\operatorname{int}(\beta')\sigma
\end{array}
\]
d'où
\[
\begin{array}{l}
(u'u)(\sigma)=u'(u(\sigma))
 =\varepsilon_2\operatorname{int}(u'(\beta))u'(\sigma)\\
\hfill{}=\varepsilon_2\varepsilon'_2
 \operatorname{int}(\deleted{\unsure{\beta'}}u'(\beta)\unsure{\beta')\sigma}
\end{array}
\]
avec \(u'(\beta)\beta'\in\widehat{\Pi}_0\), ce qui par l'unicité prouve
\[
(45)\qquad
\beta''=u'(\beta)\beta',\qquad
\varepsilon''_2=\varepsilon'_2\varepsilon_2
\]
% source-position: middle, b-coordinate multiplication and underbraces.
d'où on conclut bien
\[
\begin{array}{ccc}
b''&=&\underbrace{b'}_{\beta'^{-1}u'}\ \underbrace{b}_{\beta^{-1}u}\\
\Vert&&\\
\beta''^{-1}u''=\beta''^{-1}(u'u)&&
\end{array}
\]
\[
\begin{array}{lll}
\text{i.e.}&\beta''^{-1}u'u=\beta'^{-1}u'\beta^{-1}u
 &\text{i.e.}\quad\beta''^{-1}=\beta'^{-1}u'(\beta^{-1})\\
\text{i.e.}&\beta''=u'(\beta)\beta'
 &\text{qui n'est autre que (45).}
\end{array}
\]

% source-position: lower middle, epsilon_3 identity and (46).
La relation \(\varepsilon''_3=\varepsilon'_3\varepsilon_3\) est aussi
évidente, \unsure{p.ex.} sur la relation (qui contient \(\varepsilon_3\),
et de \(\overline{\mathrm m}\) \(\varepsilon'_3,\varepsilon''_3\)
\[
(46)\qquad
u(\rho)\deleted{\text{ est}}\ \mathcal{G}^{+\wedge}\text{-conj. à }\rho^{\varepsilon_3}.
\]
\deleted{\unsure{quand on}}

% source-position: lower third, hypothetical a-coordinate law, (*) and final displays.
Il reste donc à examiner la relation hypothétique
\[
a''=a'a
\]
qui \deleted{\unsure{s'écrit|serait}} implique
\[
(*)\qquad\delta(a'')=\delta(a')\delta(a)
\]
mais comparant (43) (43 bis) (43 ter) on trouve la relation
\[
\deleted{\unsure{\delta}}\delta(u')\delta(u)
 =\tau_\gamma^{\nu'_3}\delta(a')\tau_\gamma^{\nu_3}\delta(a)
\]
\[
\delta(u'\unsure{u''|u})=\tau_\gamma^{\nu''_3}\delta(a'')
\]
\hfill i.e.
\mspage{583}

% Printed archive footer outside the manuscript frame (not manuscript):
% 583
% https://grothendieck.umontpellier.fr
% Université de Montpellier, UMR CNRS 5149, Case courrier 051, Place Eugène Bataillon
% 34095 Montpellier cedex 5 - France

\scanpage{585}
\deleted{\(\unsure{\tau_\gamma^{\nu''_3}}\)}
\mspage{596}
% source-position: upper cancelled sentence, identities and (47).

\deleted{de sorte que la relation (*) équivaut à}
\[
\tau_\gamma^{\nu'_3+\nu_3}\delta(a'')
 =\tau_\gamma^{\nu'_3}\delta(a')\tau_\gamma^{\nu_3}\delta(a)
\]
ou encore
\[
\deleted{\tau_\gamma^{\unsure{\nu_3}}}\delta(a'')
 =\deleted{\unsure{\tau_\gamma}}\inserted{\tau_\gamma^{\nu_3}}
 \delta(a')\tau_\gamma^{\nu_3}\delta(a)
 \deleted{\unsure{\tau_\gamma^{\nu_3}}}
\]
i.e.
\[
(47)\qquad
\delta(a'')=\deleted{\unsure{\operatorname{int}}}
 \tau_\gamma^{\nu_3}(\delta(a'))\cdot\delta(a).
\]
% source-position: upper middle, (48) and its left marginal note.
Donc la relation (*) équivaut à
\[
(48)\qquad
\begin{array}{lll}
\delta(a')=\tau_\gamma^{\nu_3}(\delta(a'))
 &\text{i.e.}&[\delta(a'),\tau_\gamma^{\nu_3}]=1\quad
 \deleted{\unsure{\text{ou}}}\\
 &\text{i.e.}&[\delta(u'),\tau_\gamma^{\nu_3}]=1\quad
 \deleted{\unsure{\text{(car}}}
\end{array}
\]
\margin{puisque \(\delta(u)=\tau_\gamma^{\nu_3}\delta(a)\)}

% source-position: middle, commutation and heuristic comparison with Gamma_Q.
Cela montre que si \(\nu_3=1\) i.e. \(\varepsilon_3=-1\), alors
la relation (*) équivaut à la \uline{commutation}
\deleted{\unsure{comme}}
\uline{de} \(\delta(u')\) \deleted{\(\unsure{\delta(a')}\)}
\uline{avec} \(\tau_\gamma\),
\deleted{\unsure{commutation i.e. à}}.
C'est une condition draconienne sur un élément de \(\Gamma_{\mathbb Q}^{\sim}\)
\unsure{que|qui}, pour un élément de \(\Gamma_{\mathbb Q}\)
lui-\(\overline{\mathrm m}\) \unsure{dans sa version}, implique que
\(u'\in\{1,\tau_\gamma\}\) --- et par des arguments heuristiques, on a vu
qu'il devait \(\overline{\mathrm m}\) en être ainsi pour \unsure{tout élément}
de \(\Gamma_{\mathbb Q}^{\sim}\). Il apparait que le succès de la présentation
du \S\,48 \uline{\(\overline{\mathrm{XI}}\)} tient au fait que
\unsure{dans} \unsure{les situations|la situation}
\unsure{envisagées|envisagée} dans \unsure{les|le} cas,
le groupe \(\gamma_{\rho,\sigma}\simeq\widehat{\mathbb Z}^{*}\)
était \uline{commutatif}, \inserted{ou \(\tau_\gamma\) \uline{central}.}
% source-position: lower third, inclusion and restricted multiplicativity; cancelled left margin.
En dehors de ce cas, la quantité
\diagram{\begin{tikzcd}
a=\alpha^{-1}u\in Z_\rho^{\unsure{\sharp}}
 \arrow[r,hook] & G_{\rho,\sigma}
\end{tikzcd}}
dépend multiplicativement de l'élément
\(u\in\deleted{\unsure{Z_{\rho,\sigma}}}\mathcal M^{!}[0]\),
seulement à condition de se borner à des \(u\in\mathcal M^{!}[0]'\).

\margin{\deleted{Il suffit \(\overline{\mathrm m}\) que \(\tau_\gamma\)
soit \unsure{centralisé|central} par}
\diagram{\begin{tikzcd}
\deleted{\gamma^{!}_{\rho,\sigma}=\operatorname{Ker}(\gamma_{\rho,\sigma}}
 \arrow[r,"\deleted{\unsure{\varepsilon_3}}"'] & {}
\end{tikzcd}}}
\mspage{584}

% Printed archive footer outside the manuscript frame (not manuscript):
% 584
% https://grothendieck.umontpellier.fr
% Université de Montpellier, UMR CNRS 5149, Case courrier 051, Place Eugène Bataillon
% 34095 Montpellier cedex 5 - France

\scanpage{586}
% source-position: upper text and homomorphism diagram.

Nous allons essayer de sauver la mise, en trouvant un
\uline{homomorphisme} canonique
\diagram{\begin{tikzcd}
\mathcal M^{!}[0] \arrow[r] \arrow[d,equal]
 & \underbrace{Z_{\rho,\sigma}^{!}
 \times_{\gamma_{\rho,\sigma}}\deleted{\unsure{Z}}N_\rho^{*}
 \times_{\gamma_{\rho,\sigma}}N_\sigma^{*}}
 _{S_0\Gamma_{\rho,\sigma}^{!}\underset{\text{indice }2}{\subset}
 S_0\Gamma_{\rho,\sigma}} \\
Z_{\rho,\sigma}^{!} &
\end{tikzcd}}

\margin{\uline{NB} \unsure{on a}
\diagram{\begin{tikzcd}
S_0\Gamma_{\rho,\sigma}^{!}/L_0^{\Pi}\arrow[r,"\sim"]
 & S_0\Gamma_{\rho,\sigma}/L_0^{\mathcal G}
\end{tikzcd}}
\uuline{alors}
\diagram{\begin{tikzcd}
\gamma_{\rho,\sigma}\arrow[r,"\sim"] & \Gamma_{\mathbb Q}^{\sim}.
\end{tikzcd}}}

% source-position: upper quarter, annotated triples (49).
Notons que le deuxième \unsure{mb|ens} est formé des triplets
\diagram{\begin{tikzcd}[column sep=small,row sep=small]
(49)\quad\bigl\{\bigl(u,\arrow[d,phantom,"\in"{rotate=90}]
 & a,\arrow[d,phantom,"\in"{rotate=90}]
 & \deleted{\unsure{b}}b\bigr)\mid\arrow[d,phantom,"\in"{rotate=90}]
 & u=fb=f'a,\ \text{avec }f,f'\in\mathcal G^{+\wedge}
 \text{ convenables}\bigr\} \\
\deleted{\unsure{Z}}\unsure{\mathcal M}[0]
 & N'_\rho & N'_\sigma &
\end{tikzcd}}
\[
\text{ce qui implique}\quad
\begin{array}{l}
\varepsilon_2(u)=\varepsilon_\sigma(b)\\
\varepsilon_3(u)=\varepsilon_\rho(a)
\end{array}
\]
et pour un tel triple, on aura
% source-position: upper middle, calculations (49 bis) and definition.
\[
(49\ \mathrm{bis})\quad
\left\{\begin{array}{l}
u(\deleted{\unsure{\rho}}\inserted{\sigma})
 =\operatorname{int}(\unsure{fb})
 \deleted{\unsure{\varepsilon_2(b)|\operatorname{int}(b)}}(\sigma)
 =\operatorname{int}(f)\underbrace{\operatorname{int}(b)(\sigma)}
 _{\varepsilon_2(b)(\sigma)}
 =\deleted{\operatorname{int}(f)\unsure{\varepsilon_2}}
 \inserted{\varepsilon_2(u)\operatorname{int}(f)(\sigma)}\\
u(\rho)=\operatorname{int}(f')(\rho^{\varepsilon_3(u)})
 =\operatorname{int}(f')\varepsilon_3(\tau')(\rho)
 =\operatorname{int}(f'\varepsilon_3(\tau'))(\rho)
\end{array}\right.
\]
où on pose
\[
\varepsilon_3(\tau')\overset{\mathrm{dfn}}{=}
\left\{\begin{array}{ll}
1&\text{si }\varepsilon_3=+1\\
\unsure{\tau'}&\text{si }\varepsilon_3=-1
\end{array}\right.
\qquad
\begin{array}{l}
\text{i.e. }\varepsilon_3(\tau')=\sigma^{\nu_3}\\
\text{si }\varepsilon_3=(-1)^{\nu_3}
\end{array}
\]

Si \(u\) correspond à \(\alpha,\beta\) (cf (38)), les
% source-position: middle, equivalent conjugation formulas after (49 bis).
formules (49 bis) s'écrivent \unsure{aussi|ainsi} respectivement
\[
\left\{\begin{array}{lll}
\operatorname{int}(f)(\sigma)=\operatorname{int}(\beta)(\sigma)
 &\text{i.e.}& f=\beta\sigma^i\quad(i\in\mathbb Z/4\mathbb Z)\\
\deleted{\operatorname{int}(f')\unsure{\varepsilon_3}(\rho)=}&&\\
\operatorname{int}(f'\varepsilon_3(\tau'))(\rho)
 =\operatorname{int}(\alpha)(\rho)
 &\text{i.e.}&\underbrace{f'\varepsilon_3(\tau')=\alpha\rho^j}
 _{j\in\mathbb Z/6\mathbb Z}
\end{array}\right.
\]

Donc on trouve

% source-position: lower third, Proposition and (50).
\uline{Proposition}
Pour un \(u\in Z_{\rho,\sigma}^{!}=\mathcal M[0]
\simeq\unsure{\widetilde{\mathcal M}_{0,3}}[0]\) fixé,
les \((u,a,b)\in S_0\Gamma_{\rho,\sigma}\) au dessus de \(u\)
sont exactement les éléments \deleted{de la forme}\inserted{de}
\[
(50)\quad
\left\{(u,a,b)\left|\begin{array}{lll}
a=f'^{-1}u&\inserted{\text{avec}}\quad
 f'=\alpha\rho^j\varepsilon_3(\tau')&(j\in\mathbb Z/6\mathbb Z)\\
b=f^{-1}u&\text{avec }f=\beta\sigma^i&(i\in\mathbb Z/4\mathbb Z)
\end{array}\right.\right\}
\]
\margin{où \(\alpha,\beta,\varepsilon_2,\varepsilon_3\) sont
définis pour \(u\) \unsure{via} (38)}

\deleted{On aura}\inserted{Considérons} alors l'élément
% source-position: final paragraph following (50), ending at the bottom of the scan.
\inserted{\((u,a,b)\)} au dessus de \(u\)
ayant l'expression la plus simple en termes de \(\alpha,\beta\),
en prenant \(i=\unsure{0|j},\ j=0\), correspondant donc aux choix :
\mspage{585}

% Printed archive footer outside the manuscript frame (not manuscript):
% 585
% https://grothendieck.umontpellier.fr
% Université de Montpellier, UMR CNRS 5149, Case courrier 051, Place Eugène Bataillon
% 34095 Montpellier cedex 5 - France
